% !TEX program = xelatex
% !TEX encoding = UTF-8
\documentclass[a4paper, 11pt]{ctexart}
% 页面与边距
\usepackage[margin=2cm]{geometry}
% 数学与工具宏包
\usepackage{amsmath, amsfonts, amssymb}
\usepackage{graphicx}
\usepackage{float}
\usepackage{xcolor}
\usepackage{fancyhdr}
\usepackage{etoolbox} % 用于逻辑判断
\usepackage{hyperref}
\usepackage{caption}
% ================= 字体设置 =================
% 引入 Inconsolata 等宽字体宏包，使代码英文字符显示为标准等宽字体
% \usepackage{zi4} 
% 如果你想使用系统自带的其他等宽字体（如 Consolas），
% 可以注释掉上一行 \usepackage{zi4}，并取消下面两行的注释（需使用 XeLaTeX）：
% \usepackage{fontspec}
\setmonofont{Consolas}
% ================= 页眉页脚 =================
\pagestyle{fancy}
\fancyhf{}
\cfoot{\thepage}
\renewcommand{\headrulewidth}{0pt}
% ================= 题目环境 (支持自动/自定义) =================
% ================= 模式切换开关 =================
\newif\ifshowproblem
\showproblemtrue  % 模式 1：显示题目内容 (默认状态)
% \showproblemfalse % 模式 2：隐藏题目内容 (取消这一行的注释即可切换)
% ================= 题目环境 (支持模式切换与二级标题) =================
\newcounter{problem}[section]
\renewcommand{\theproblem}{\thesection.\arabic{problem}}
\NewDocumentEnvironment{problem}{o +b}{%
    \IfNoValueTF{#1}{%
        \refstepcounter{problem}%
        \subsection*{题目 \theproblem}%
        \addcontentsline{toc}{subsection}{题目 \theproblem}%
    }{%
        \subsection*{题目 #1}%
        \addcontentsline{toc}{subsection}{题目 #1}%
    }%
    % 判断模式：如果开启了显示题目，则输出被吸收的题目内容（#2）
    \ifshowproblem
        #2%
    \fi
}{%
    % 只有在显示题目的模式下，才在底部增加间距
    \ifshowproblem
        \vspace{1ex}%
    \fi
}
% ================= 答案框 =================
\usepackage[many]{tcolorbox}
\newtcolorbox{answer}{
    enhanced,
    breakable,
    title=解答,
    colback=white,
    colframe=black!80,
    fonttitle=\bfseries,
    attach boxed title to top left={xshift=5mm,yshift=-2mm},
    boxed title style={colback=black!80, size=small},
    sharp corners,
    top=4mm, bottom=4mm, left=4mm, right=4mm,
    before skip=10pt, after skip=15pt,
    pad at break=2mm,
    bottomrule at break=1mm,
    toprule at break=1mm,
    % --- 新增以下 overlay 代码来手动绘制跨页处的边框 ---
    % 在第一页断页处的底部画线
    overlay first={\draw[black!80, line width=0.5mm] (frame.south west) -- (frame.south east);},
    % 在第二页接续处的顶部画线
    overlay last={\draw[black!80, line width=0.5mm] (frame.north west) -- (frame.north east);},
    % 如果盒子内容极长跨了三页及以上，为中间页的上下都画线
    overlay middle={\draw[black!80, line width=0.5mm] (frame.south west) -- (frame.south east); \draw[black!80, line width=1mm] (frame.north west) -- (frame.north east);}
}
% ================= GitHub 风格代码高亮 =================
\usepackage{listings}

% GitHub 风格色彩定义
\definecolor{gh_comment}{RGB}{106, 115, 125}   % 灰色注释
\definecolor{gh_keyword}{RGB}{215, 58, 73}   % 红色关键字
\definecolor{gh_string}{RGB}{3, 47, 98}      % 深蓝字符串
\definecolor{gh_constant}{RGB}{0, 92, 197}   % 蓝色常量/数字
\definecolor{gh_bg}{RGB}{246, 248, 250}      % 浅灰背景

\lstset{
    language=Python,
    backgroundcolor=\color{gh_bg},
    basicstyle=\ttfamily\small,      % 使用等宽字体 monomaple
    breaklines=true,
    showstringspaces=false,
    commentstyle=\color{gh_comment},
    keywordstyle=\color{gh_keyword}\bfseries,
    stringstyle=\color{gh_string},
    numberstyle=\tiny\color{gh_comment},
    numbers=left,
    frame=single,
    rulecolor=\color{black!10},
    xleftmargin=1em,
    tabsize=4
}
% ================= 包装代码环境 =================
\usepackage{ifthen}
% 定义 \begin{code}[语言名称] ... \end{code}
\lstnewenvironment{code}[1][]
{
    % 在环境开始时，将方括号内传入的参数（如 Python）赋给 language
    \ifthenelse{\equal{#1}{}}{}{\lstset{language=#1}}
}
{
    % 环境结束时的操作（保持为空）
}
\everymath{\displaystyle}
% ---------------- 标题 ----------------
\title{\textbf{课程作业1}}
\author{Jiuxiao}
\date{\today}

\begin{document}

\maketitle

\section{基础章节}

\begin{problem}
这里是题目1。
\end{problem}

\begin{answer}

可以使用 \textbackslash{}code 环境来插入代码：

\begin{code}[Python]
def hello_fun():
    print("Hello world")

hello_fun()
\end{code}

\end{answer}

\begin{problem}[1.1.1]
这是通过参数自定义的编号。即使手动指定了编号，下一次使用自动编号时，计数器依然会从上一题的基础位移。
\end{problem}

\begin{answer}
准确率（Accuracy）与错误率（Error Rate）的计算公式如下：

\[
\text{Accuracy} = \frac{TP + TN}{P + N}, \quad \text{Error Rate} = 1 - \text{Accuracy}
\]

\end{answer}

\begin{problem}
这是紧接在自定义编号后的自动编号。
\end{problem}
\begin{answer}
    显然。$\square$
\end{answer}

\section{附加题}
\begin{problem}
这是第二部分的题目。
\end{problem}
\begin{answer}
    文档开启了为所有数学公式启用 \textbackslash{}displaystyle参数，
    这会使得\\行内公式 $ \text{Attention}(Q,K,V)=\text{softmax}\left(\frac{QK^T}{\sqrt{d_k}}\right)V $ 和 行间公式
    \[
        \text{Attention}(Q,K,V)=\text{softmax}\left(\frac{QK^T}{\sqrt{d_k}}\right)V
    \]
    一样大，注释掉前面的 \textbackslash{}everymath\{\textbackslash{}displaystyle\} 即可恢复默认
\end{answer}

\begin{problem}
文档提供了输出纯答案的功能开关，在导言区使用\textbackslash{}showproblemtrue 为显示题目内容 (默认状态)，需要提交作业的时候使用 \textbackslash{}showproblemfalse 覆盖为隐藏题目内容模式，可以输出无题目仅答案的版本。
\end{problem}
\begin{answer}
    现在，你可以删除正文部分并且开始使用这份模板了。
\end{answer}
% \begin{answer}
%     大框测试。
%     \[
%         \text{Attention}(Q,K,V)=\text{softmax}\left(\frac{QK^T}{\sqrt{d_k}}\right)V
%     \]\[
%         \text{Attention}(Q,K,V)=\text{softmax}\left(\frac{QK^T}{\sqrt{d_k}}\right)V
%     \]\[
%         \text{Attention}(Q,K,V)=\text{softmax}\left(\frac{QK^T}{\sqrt{d_k}}\right)V
%     \]
%     \[
%         \text{Attention}(Q,K,V)=\text{softmax}\left(\frac{QK^T}{\sqrt{d_k}}\right)V
%     \]\[
%         \text{Attention}(Q,K,V)=\text{softmax}\left(\frac{QK^T}{\sqrt{d_k}}\right)V
%     \]\[
%         \text{Attention}(Q,K,V)=\text{softmax}\left(\frac{QK^T}{\sqrt{d_k}}\right)V
%     \]\[
%         \text{Attention}(Q,K,V)=\text{softmax}\left(\frac{QK^T}{\sqrt{d_k}}\right)V
%     \]\[
%         \text{Attention}(Q,K,V)=\text{softmax}\left(\frac{QK^T}{\sqrt{d_k}}\right)V
%     \]\[
%         \text{Attention}(Q,K,V)=\text{softmax}\left(\frac{QK^T}{\sqrt{d_k}}\right)V
%     \]\[
%         \text{Attention}(Q,K,V)=\text{softmax}\left(\frac{QK^T}{\sqrt{d_k}}\right)V
%     \]\[
%         \text{Attention}(Q,K,V)=\text{softmax}\left(\frac{QK^T}{\sqrt{d_k}}\right)V
%     \]\[
%         \text{Attention}(Q,K,V)=\text{softmax}\left(\frac{QK^T}{\sqrt{d_k}}\right)V
%     \]\[
%         \text{Attention}(Q,K,V)=\text{softmax}\left(\frac{QK^T}{\sqrt{d_k}}\right)V
%     \]
% \end{answer}
\end{document}